\documentclass[letterpaper]{article}
\usepackage[submission]{aaai2026}
\usepackage{xeCJK}
\setCJKmainfont{SimSun}
\setCJKsansfont{SimHei}
\setCJKmonofont{FangSong}
\usepackage{times}
\usepackage{helvet}
\usepackage{courier}
\usepackage[hyphens]{url}
\usepackage{graphicx}
\urlstyle{rm}
\def\UrlFont{\rm}
\usepackage{natbib}
\usepackage{caption}
\frenchspacing
\setlength{\pdfpagewidth}{8.5in}
\setlength{\pdfpageheight}{11in}
\usepackage{amsmath}
\usepackage{amssymb}
\usepackage{booktabs}
\usepackage{multirow}
\usepackage{tikz}
\usepackage{placeins}
\usetikzlibrary{arrows.meta,positioning,calc,fit}
\pdfinfo{
/TemplateVersion (2026.1)
}
\setcounter{secnumdepth}{0}

\title{潜变量桥接匹配：基于最优传输耦合流匹配的高效域级艺术风格迁移}
\author{匿名投稿}
\affiliations{}
\begin{document}
\maketitle

\begin{abstract}
本文研究域级艺术风格迁移中风格强度、内容保留与计算成本三者之间的实际约束。我们提出潜变量桥接匹配（LBM），一种高效的潜空间框架，融合了小批量最优传输（OT）耦合、时间条件流匹配、语义路由与终端切片 Wasserstein（SWD）正则化。OT耦合为非配对监督提供合理的风格域端点，流匹配学习稳定的潜变量传输路径，终端SWD将积分后的输出对齐到目标域的块统计特征。在严格的750输出协议（5域×5风格×30图像）下，我们的3.9M参数模型在最接近的高效基线SaMST上达到近乎一致的CLIP风格（0.716对0.719），提升LPIPS和EC，训练成本降低约21.8倍，推理速度相当，同时在核心迁移指标上优于更重的S2WAT基线。这些结果将LBM置于域级艺术风格迁移的有利效率-质量帕累托前沿上。
\end{abstract}

\section{引言}
艺术风格迁移长期以来是分离图像内容与视觉外观的试验床。经典神经风格迁移展示了卷积特征和Gram统计量可合成分令人信服的艺术纹理\cite{gatys2016image}，而前馈式任意风格迁移通过归一化风格调制实现了实时部署\cite{huang2017adain}。近年来，注意力机制、Transformer、免训练和扩散驱动的方法进一步提升了灵活性和风格保真度\cite{park2019sanet,deng2022stytr2,hong2023aespa,huang2024aesfa,gim2024cast,chung2024styleid}。然而，实用的多风格艺术迁移仍然暴露于一个持续的三角矛盾中：强目标风格表达、忠实内容保留与低计算成本三者难以兼得。

这种矛盾在域级多风格迁移中尤为显著，其目标不是单参考图像模仿，而是跨多个目标域的高效摊销迁移。在我们的严格750输出协议下，StyleID获得最高的原始CLIP风格，但LPIPS达到0.750，表明严重的内容漂移；SaMST实现了强风格和内容得分，但产生可见的模糊和颗粒状纹理伪影；S2WAT作为更重的Transformer基线，LPIPS明显更高（0.526）。这些观察表明，仅凭风格/内容指标是不够的——一个有竞争力的系统还必须在感知质量和纹理真实性上接受评估。

我们认为域级艺术风格迁移更适合视为三个耦合的子问题：为非配对监督选择合理的风格域端点、学习通向该端点的稳定传输路径、以及确保实际积分后的输出落在目标域分布中。LBM分别用最优传输（OT）耦合、时间条件潜变量流匹配和终端SWD正则化来解决这三个问题，而语义交叉注意力在紧凑的潜空间中提供轻量级局部风格路由。

实验上，LBM保持了最强高效基线SaMST的质量水平，同时将训练成本降低了超过20倍，推理速度相当。在更重的S2WAT基线上，LBM不仅在效率上更优，在核心风格-内容权衡指标上也更强。

\subsection{贡献}
\begin{itemize}
    \item \textbf{一种域级艺术风格迁移的结构化潜变量传输形式化}，将非配对监督分解为基于OT的端点构建、路径流学习和终端端点正则化。
    \item \textbf{一种带有语义路由和学习型空间风格先验的轻量级潜变量主干网络}，无需逐图像优化、扩散采样或参考特定编码即可实现高效域级迁移。
    \item \textbf{在最近的高效基线SaMST上实现了强效率-质量结果}：CLIP风格近乎一致（0.716对0.719）、LPIPS改善（0.451对0.466）、EC更高（0.393对0.384）、训练成本降低约22倍。
    \item \textbf{在更重的S2WAT基线上取得持续优势}：CLIP风格更高、LPIPS显著更低、EC更高、模型小16.7倍——证明效率提升并非以牺牲迁移质量为代价。
\end{itemize}

\section{相关工作}
\paragraph{经典与任意神经风格迁移。}
神经风格迁移始于基于优化的特征重构和Gram矩阵风格损失\cite{gatys2016image}。AdaIN通过对齐通道级特征统计实现了实时任意迁移\cite{huang2017adain}。基于注意力的方法如SANet通过风格注意力特征聚合改进了局部风格匹配\cite{park2019sanet}。近期审美和模式感知方法如AesPA-Net和AesFA进一步强调模式可重复性、审美特征和高质量任意风格化\cite{hong2023aespa,huang2024aesfa}。这些方法展示了特征统计和局部匹配的价值，但它们通常不显式建模包含终端分布约束的风格条件传输路径。

\paragraph{Transformer与多风格迁移。}
基于Transformer的风格迁移方法如StyTr2和S2WAT利用注意力机制改善内容-风格交互和长程特征建模\cite{deng2022stytr2,xia2024s2wat}。多风格方法将许多风格摊销在一个紧凑系统中；SaMST是最新的代表，具有可插拔风格表示、低参数量和快速推理\cite{liu2024samst}。这些系统与我们高效多风格迁移的目标直接相关。与纯表示风格注入不同，我们的方法将语义交叉注意力置于潜桥接目标内，因此终端分布匹配和动能正则化显式塑造了风格-内容权衡。

\paragraph{免训练与扩散方法。}
免训练和扩散驱动的方法提供了通向灵活风格迁移的另一条路径。CAST使用VQ自编码器实现内容自适应的免训练风格迁移\cite{gim2024cast}，而StyleID将风格注入大规模扩散模型而无需训练\cite{chung2024styleid}。这些方法可以在视觉上很强，但扩散采样和语义漂移可能限制其在快速域级迁移中的适用性。我们的方法瞄准一个互补的效率区间：紧致潜积分与风格ID条件，而非迭代扩散采样或逐风格优化。

\section{方法}
\subsection{概述}
LBM最好被理解为一种用于域级艺术风格迁移的结构化潜变量传输设计。设$z_0 \in \mathbb{R}^{4 \times 32 \times 32}$为内容图像的VAE潜变量，$s$为目标风格域标识。OT耦合选择合理的风格域端点，流匹配学习通向它们的传输路径，终端SWD矫正积分轨迹的落点，语义路由控制风格在局部写入的位置。

\subsection{非配对潜域采样与OT耦合}
我们直接在预计算的VAE潜变量上进行无配对监督训练。每个小批量从所有源域采样内容潜变量$\{z_c\}$，从请求的目标域采样目标风格潜变量$\{z_s\}$。通过均匀目标域采样自然包含恒等对，这稳定了同域映射上的学习传输。

我们使用切片Wasserstein成本构建小批量传输耦合，而非随机目标分配。对于一批内容潜变量$\{z_c^{(i)}\}$和目标风格潜变量$\{z_s^{(j)}\}$，SWD成本矩阵为
\begin{equation}
    C_{ij} = \mathrm{SWD}\bigl(z_c^{(i)}, z_s^{(j)}\bigr) .
\end{equation}
然后通过Sinkhorn算法获得熵最优传输计划：
\begin{equation}
    \pi^\star = \arg\min_\pi \sum_{ij} C_{ij}\pi_{ij} + \epsilon H(\pi),
\end{equation}
其中$H(\pi)$是熵正则化项。每个内容潜变量被匹配到其耦合目标$\tilde{z}_1 \sim \pi^\star(z_s|z_c)$。

\subsection{时间条件潜变量流匹配}
核心是一个时间条件速度场$v_\theta(z_t,t,s)$，沿随机桥传输潜变量。在随机采样的时间$t\sim\mathcal{U}(0,1)$，桥状态为
\begin{equation}
    z_t = (1-t)z_0 + t\tilde{z}_1 + \sigma\sqrt{t(1-t)}\,\epsilon,
    \quad \epsilon\sim\mathcal{N}(0,I),
\end{equation}
目标速度为传输的导数：
\begin{equation}
    u_t = \tilde{z}_1 - z_0 + \sigma\frac{1-2t}{2\sqrt{t(1-t)}}\,\epsilon .
\end{equation}
流匹配目标监督网络预测此速度：
\begin{equation}
    \mathcal{L}_{\mathrm{FM}} = \mathbb{E}_{z_0,\tilde{z}_1,t,\epsilon}\,
    \bigl\|v_\theta(z_t,t,s) - u_t\bigr\|^2 .
\end{equation}

\paragraph{架构。}
时间条件主干使用正弦时间嵌入，加到风格编码上。实现类似UNet的潜变量主干，在$16\times16$分辨率有四个语义交叉注意力体块。目标风格通过可学习的空间先验$P_s\in\mathbb{R}^{128\times16\times16}$表示。语义交叉注意力将风格先验注入内容特征，跳跃连接保留空间结构。

\subsection{终端SWD正则化}
除流匹配监督外，终端损失将欧拉积分后的输出正则化到目标风格分布。设$\Phi_\theta(z_0,s)$为从$z_0$用风格$s$积分$v_\theta$得到的端点。终端损失为
\begin{equation}
    \mathcal{L}_{\mathrm{term}} = \mathrm{SWD}\bigl(\Phi_\theta(z_0,s),\,\tilde{z}_1\bigr),
\end{equation}
其中SWD用64个投影和块大小$\{3,5,7,15\}$比较生成的端点块与耦合目标块。

完整训练目标结合流匹配、终端分布匹配和动能路径正则化：
\begin{equation}
    \mathcal{L} = \mathcal{L}_{\mathrm{FM}} + \lambda_{\mathrm{term}}\mathcal{L}_{\mathrm{term}} + \lambda_{\mathrm{kin}}\mathcal{L}_{\mathrm{kin}},
\end{equation}
其中$\mathcal{L}_{\mathrm{kin}} = \mathbb{E}\|v_\theta\|_2^2$惩罚过大的潜变量位移。主要配置使用$\lambda_{\mathrm{term}}=20.0$和$\lambda_{\mathrm{kin}}=1.0$。

\paragraph{互补角色。}
三个组件解决不同的子问题。OT耦合单独提供合理端点但不提供传输动力学。流匹配学习传输路径但不保证积分端点匹配目标域纹理统计。终端SWD矫正这种端点失配，而语义交叉注意力使风格传输在局部上有意义。这种分离是有意为之：弱化OT耦合消除端点引导；移除终端SWD使端点风格崩溃；移除流匹配路径监督破坏传输稳定性。

\paragraph{有限的理论声明。}
我们的方法不是严格的薛定谔桥求解器。它不估计正向和反向随机得分，也不优化完整的熵对偶目标。桥视角被用作一个可操作的训练原则：流匹配路径约束、终端分布正则化器和动能惩罚。这种有限解释也在经验上得到反映——移除终端约束弱化风格，而移除路径惩罚破坏内容。

\section{实验}
\subsection{协议与指标}
我们在严格的750输出协议下评估：5个源域乘以5个目标域乘以30张源图像。五个域为photo、Hayao、Monet、Van Gogh和Cezanne。我们报告CLIP风格、CLIP内容、LPIPS内容和复合得分
\begin{equation}
    \mathrm{EC} = \mathrm{CLIP\text{-}style}\cdot(1-\mathrm{LPIPS}).
\end{equation}
为进行伪影敏感分析，我们还报告MUSIQ、MANIQA、DISTS内容、高频块KID、FFT斜率误差和浅层Gram微观统计量。我们对Ours、SaMST、S2WAT、StyleID、AdaIN和CAST执行端到端的严格750协议。

\subsection{主要比较}
\begin{table*}[t]
\centering
\caption{严格750核心对比与效率指标。$\uparrow$ = 越高越好，$\downarrow$ = 越低越好。推理时间为750张图像的端到端测量。$\dagger$标记的训练时间由外推获得。}
\label{tab:main}
\small
\setlength{\tabcolsep}{5pt}
\resizebox{\textwidth}{!}{%
\begin{tabular}{@{}lrrrrrrr@{}}
\toprule
方法 & 参数量 & CLIP-S$\uparrow$ & LPIPS$\downarrow$ & EC$\uparrow$ & 推理 & 秒/张 & 训练时间 \\
\midrule
Ours ep7 & 3.9 M & 0.716 & 0.451 & 0.393 & 40.0 s & 0.053 & 309.9 \\
SaMST & 6.0 M & 0.719 & 0.466 & 0.384 & 39.8 s & 0.053 & 6768.7$\dagger$ \\
S2WAT & 65 M & 0.714 & 0.526 & 0.338 & 45.1 s & 0.060 & 10600$\dagger$ \\
AdaIN v32k & $\sim$5 M & 0.713 & 0.630 & 0.264 & 9.3 s & 0.012 & 9220.4 \\
\midrule
StyleID & SD基础 & 0.760 & 0.750 & 0.190 & 603.3 s & 0.804 & 免训练 \\
CAST & 7.0 M & 0.665 & 0.726 & 0.182 & 75.5 s & 0.101 & 免训练 \\
\bottomrule
\end{tabular}
}
\smallskip
\footnoterule\small
$\dagger$外推值。
\end{table*}

关键对比是对SaMST和S2WAT。最强的高效基线是SaMST。相对于SaMST，我们的方法在CLIP风格上基本持平（0.716对0.719），LPIPS改善（0.451对0.466），EC更高（0.393对0.384），参数量从6.0M减至3.9M，训练成本降低约22倍，推理速度相当。这是本文的核心实验发现：LBM保持了最强高效基线的质量水平，同时大幅降低了训练成本。

相对于S2WAT，LBM不仅更小、训练成本更低，而且在核心迁移指标上也更强：CLIP风格更高、LPIPS显著更低、EC更高。这一对比表明效率提升不是以牺牲迁移质量为代价的——LBM在更重基线上呈现帕累托优势。

AdaIN推理速度仍然更快，但明显低于高质量前沿（LPIPS 0.630、EC 0.264），而StyleID达到最高原始CLIP风格（0.760）的代价是严重的内容漂移（LPIPS 0.750、CLIP内容0.552）和慢得多的推理（750张图603秒）。

\subsection{对SaMST的伪影敏感对比}
由于SaMST在核心指标下是最近的高效基线，我们将其作为伪影聚焦诊断的主要目标。尽管其标题分数具有竞争力，放大检查揭示了标准风格/内容指标未能充分惩罚的模糊和颗粒状局部纹理。SSIM和边缘重叠等指标奖励低频结构保留；因此，即使局部纹理视觉上不干净，只要物体边界可见，颗粒状输出也能获得高分。

为捕捉这种失败模式，我们检查表\ref{tab:artifact}中的伪影敏感诊断指标。Ours将MUSIQ从36.10提升至49.21，MANIQA从0.314提升至0.406，DISTS内容从0.294降至0.248，HF-块KID从6.76降至4.17，FFT斜率误差从1.054降至0.547。这些改进表明我们的方法虽然在原始风格上略微保守，但在块级别上视觉上更干净——与定性网格图一致。

\begin{table}[t]
\centering
\caption{Ours与SaMST的伪影敏感对比。}
\label{tab:artifact}
\begin{tabular}{lrr}
\toprule
指标 & Ours e7 & SaMST \\
\midrule
MUSIQ $\uparrow$ & 49.2059 & 36.0950 \\
MANIQA $\uparrow$ & 0.4057 & 0.3139 \\
DISTS内容 $\downarrow$ & 0.2477 & 0.2943 \\
HF-块KID $\downarrow$ & 4.1694 & 6.7598 \\
FFT斜率误差 $\downarrow$ & 0.5473 & 1.0536 \\
Gram微观 $\downarrow$ & 0.0798 & 0.0947 \\
\bottomrule
\end{tabular}
\end{table}

\subsection{破坏性消融实验}
\begin{table}[t]
\centering
\caption{严格750协议下的破坏性消融实验。}
\label{tab:ablation}
\resizebox{\columnwidth}{!}{%
\begin{tabular}{@{}lrrrr@{}}
\toprule
变体 & CLIP-S & CLIP-C & LPIPS & EC \\
\midrule
D0 full & 0.7014 & 0.8022 & 0.4593 & 0.3791 \\
D1 无终端SWD & 0.6708 & 0.8989 & 0.3490 & 0.4368 \\
D2 无动能 & 0.7159 & 0.6624 & 0.6375 & 0.2596 \\
D8 强色彩损失 & 0.6923 & 0.6629 & 0.5675 & 0.2994 \\
\bottomrule
\end{tabular}
}
\end{table}

消融结果暴露了模型机制。移除终端SWD在内容保留上出色但风格弱，表明终端分布匹配是主要的风格驱动因素。移除动能正则化保留了风格但内容崩溃，确认了路径能量惩罚的必要性。强色彩损失同时损害内容和分布质量，因此主线中不使用简单色彩匹配。

\subsection{权重扫描与理论开关}
40轮类别权重扫描训练了20个采样配方，各$K\in\{1,2\}$，8个epoch，共评估320个检查点。最佳EC结果为K2均衡默认第3epoch：CLIP风格0.6980、CLIP内容0.8727、LPIPS 0.3777、EC 0.4343。扫描显示更复杂的类别权重并不持续优于均衡采样。

我们还评估了可选的理路开关。Sinkhorn归一化语义路由和熵门控动能正则化改善了EC/内容但降低了原始CLIP风格。色彩传输略微提升原始风格但损害了LPIPS和内容。

\subsection{计算资源与可复现性}
端到端推理和训练时间列在表\ref{tab:main}的主指标旁。所有计时在NVIDIA RTX 4070 Laptop GPU（8 GB显存）上使用相同的严格750协议测量。$\dagger$标记的训练时间来自单epoch探针的外推。StyleID和CAST本质上是免训练的。

所有严格750评估使用相同的域协议和预处理流程。我们还评估了40轮扫描中的每个epoch，共320个检查点评估。为提高可复现性，每次训练运行保存了解析后的运行时配置、源文件快照、检查点和逐epoch CSV日志。Python、NumPy、PyTorch和DataLoader工作线程种子在每次运行中固定。

\section{讨论与局限性}
我们不声称在所有方法上最大化原始风格识别；相反，我们声称在高效域级体制中实现了显著更强的效率-质量权衡。核心发现是LBM在CLIP风格上保持与SaMST近乎一致（0.716对0.719），提升LPIPS和EC，训练成本降低约22倍，并在核心迁移指标上优于更重的S2WAT基线——这是一个帕累托占优的效率-质量操作点。

当前研究仍有几个局限性。首先，风格表示是域级而非参考图像特定的；这使得推理高效但限制了对独特单图像笔触模式的复现能力。第二，欧拉积分故意保持简单，因此当潜路径外推到更激进的设置或更高分辨率时，离散化误差可能累积。第三，尽管核心指标和效率分析对比较的基线是完整的，完整伪影诊断包仅对方法子集运行，尚未收集用户研究数据。

\section{结论}
我们提出了基于最优传输耦合潜变量流匹配的高效域级艺术风格迁移方法。该方法结合三个互补目标：流匹配损失在随机潜桥时间上监督时间条件速度场；终端切片Wasserstein匹配将积分端点拉向目标风格分布；动能正则化约束路径位移以保留内容。在严格的750输出评估下，结果3.9M参数模型达到了具有竞争力的CLIP风格（0.716）、高风格方法中最低的LPIPS（0.451）、最佳EC得分（0.393）以及比更重前馈基线显著更低的训练成本——一个域级艺术风格迁移的帕累托高效操作点。

\FloatBarrier
\bibliography{refs}

\clearpage
\small
\setlength{\itemsep}{1pt}
\setlength{\parskip}{2pt}
\section*{可复现性检查表}
...（检查表部分翻译略）

\end{document}
